\section{Pretraining and financial proximity}

Fix the dated context. Let $F(h)$ be the ideal task prediction for an
issuer history $h$, and let $s(h)$ be its economic target.

\begin{proposition}[Task alignment]
  \label{prop:alignment}
  For a nonempty target space, every representation from which $F$ can
  be recovered also admits an unrestricted readout of $s$ if and only if
  $F(h)=F(h')$ implies $s(h)=s(h')$ for every pair of histories.
  \lean{FreeFoundations.Alignment.optimal_imp_sufficient_iff}
\end{proposition}

To see necessity, take the representation to be $F$ itself. For
sufficiency, assign each value of $F$ the common exposure of the histories
giving that prediction and extend the readout arbitrarily off the image.
A particular encoder can retain distinctions beyond those the task requires.

Task diversity can also make recovery stable. A task measuring the sum
of financing and operating exposures gives the same prediction for
$(1,0)$ and $(0,1)$. Measuring their difference as well distinguishes
the two directions. More generally, suppose a vector of task predictions
satisfies $\tau=Ca+v$ for risk-scaled exposure $a$ and mismatch $v$.
If $\norm{Cu}\geq\kappa\norm u$ in every direction for $\kappa>0$ in
finite dimension, there is a linear decoder $L$ satisfying
\[
  \norm{L\widehat\tau-a}\leq
  (\norm{\widehat\tau-\tau}+\norm v)/\kappa.
\]
The margin describes the economic directions covered by the tasks.
The issuer study varies the financial fields supplied to the encoder.
\lean{FreeFoundations.Recovery.exists_decoder}
\lean{FreeFoundations.Recovery.exposure_error_mismatch}

\begin{proposition}[Recoverability and geometric identification]
  \label{prop:recovery}
  Let $s:H\to S$ and $z:H\to Z$, with $S$ nonempty. An unrestricted
  pointwise readout $r:Z\to S$ with $s=r\circ z$ exists exactly when
  $s$ is constant on the fibres of $z$. For any bijection $g:Z\to Z$,
  the recoding $g\circ z$ recovers the same target with readout
  $r\circ g^{-1}$. Such a recoding can change nearest neighbours.
\end{proposition}

Consider $A=(1,1)$, $B=(1,0)$, $C=(1,4)$ and $D=(1,5)$.
Stretch the first coordinate by four. A readout of the original vectors
can undo the stretch. For peer selection, normalize each vector to unit
length and compare chord distances. The closest peers of $A,B,C,D$ are
initially $C,A,D,C$ and become $B,A,D,C$. Joining firms whenever either
selects the other therefore removes the $AC$ edge and changes peer weights.
\lean{FreeFoundations.Neighbours.symmetrized_chord_edge_flip}

For ten neighbours, a sufficient certificate uses distances between
estimated exposures. If every issuer's exposure error is bounded by
$\delta_j\leq\delta_{\max}$, a strict gap above
$2(\delta_i+\delta_{\max})$ between issuer $i$'s tenth and eleventh
estimated distances preserves its selected set. Applying this certificate
to native hidden-state distances would require evidence relating them
to exposure distances. Appendix~\ref{app:proofs} gives the derivation.

\paragraph{Distances and weights.}
For the US benchmark, we compute Spearman chord distance for each
characteristic over paired annual observations. Pairs with fewer than
two common years or an undefined correlation receive the neutral
distance $\sqrt{2}$. The joint distance is the root mean square
across the four characteristics. We set the diagonal of $K$ to zero,
set $K_{ij}=\exp(-d_{ij})$ off the diagonal, and divide each row by
its sum. On normalized nonzero embeddings the chord distance is the
Euclidean distance between unit vectors. To construct a binary
$k$-nearest-neighbour graph we join two firms if either selects the
other before row normalization. Both constructions use a symmetric raw affinity.

\paragraph{Spatial specification.}
Write $y_i$ for median quarterly excess return and $b_i$ for the
issuer's estimated CAPM beta. We compare the conventional
cross-section to the spatial lag model
$y=\alpha\mathbf{1}+\gamma b+\rho Wy+u$, with $\alpha$ either fixed
to zero or estimated. The corresponding spatial-error specification
applies the multiplier to the error and leaves the systematic term
unchanged. With $\alpha=0$, the peer average may absorb the
cross-sectional mean. With an intercept, exact uniform off-diagonal
weights make the added regressor a linear combination of a constant
and $y$, so the spatial term requires variation beyond the sample
mean and the mechanical own-observation correction.

Peer weights also enter the spatial multiplier.
Proposition~\ref{prop:multiplier}
bounds the resulting change in implied covariance for a common roster,
spatial coefficient and innovation covariance.

\paragraph{Estimation and uncertainty.}
For the characteristic benchmark the estimated quantities include the
market betas, the cross-sectional premium, the intercept when
included, the spatial coefficient, and the error variance. For the
encoder comparisons, we fix the model weights, the rule for
selecting historical examples, the preprocessing and the distance measure.
To assess variation across quarters at the estimated loadings and spatial
parameters, we resample the paired likelihood contributions in common time
blocks. In the fixed-cohort exercise, where issuer betas and median returns
enter a single cross-section, we recompute both quantities in each refit.
